\documentclass[11pt,a4paper]{article}
% Shared preamble for RuPhO English problem statements (match T1 2026 style).
% Use: \input{latex/rupho_en_preamble.tex} after \documentclass.
\usepackage[margin=1in]{geometry}
\usepackage{amsmath,amssymb}
\usepackage{graphicx}
\usepackage{float}
\usepackage{enumitem}
\usepackage{microtype}
\usepackage{caption}
\usepackage{tabularx}
\usepackage{hyperref}

\captionsetup{font=small,labelfont=bf,skip=6pt}

\setlist[itemize]{leftmargin=1.2cm}
\setlist[enumerate]{leftmargin=1.2cm}

% One outer box per subproblem: [id | statement | points] (no inner rules)
\newcommand{\problembox}[3]{%
  \par\addvspace{0.42em}%
  \noindent
  \setlength{\fboxsep}{7.5pt}%
  \setlength{\fboxrule}{0.95pt}%
  \fbox{%
    \begin{minipage}{\dimexpr\linewidth-2\fboxsep-2\fboxrule\relax}
    \begin{tabularx}{\linewidth}{@{}>{\raggedright\arraybackslash\bfseries}p{2.1em} @{\hspace{0.9em}} >{\raggedright\arraybackslash}X @{\hspace{0.9em}} >{\raggedleft\arraybackslash\bfseries}p{2.4em}@{}}
    #1 & #2 & #3
    \end{tabularx}
    \end{minipage}%
  }%
  \par\addvspace{0.32em}%
}

\title{\textbf{Phase transition modeling}}
\author{\normalsize Y21 (2021) --- English translation}
\date{}
\begin{document}
\maketitle

Hint: The running time for this task can be quite long for some parameter values. Therefore, you should choose the parameters and operating time so that, on the one hand, you achieve reasonable accuracy, and on the other hand, have time to make the necessary measurements. To speed up calculations, you can run several copies of the same program in parallel with different parameters. Errors should only be assessed at points where this is explicitly requested.

The two-dimensional Ising model is one of the simplest models of the paramagnetic-ferromagnetic phase transition. The model is a square lattice, at the nodes of which magnetic moments are located. Each of these magnetic moments can be directed in one of two directions - up or down. The state of the magnetic moment is given by the variable $s_i$, where the index $i$ denotes a lattice site. $s_i = 1$ corresponds to a magnetic moment directed upward, and $s_i =-1$ – downward. The energy of such a system is determined by the interaction of nearby magnetic moments. If two magnetic moments are co-directed, the contribution of this pair to the energy is $-J$, if they are oppositely directed, then $+J$. Here $J$ is a certain constant with the dimension of energy, characterizing the strength of interaction between magnetic moments. If the system is placed in an external magnetic field, each magnetic moment will have an additional energy $h$ if it is directed against the field, and an energy $-h$ if it is directed along the field. Value $h = \mu B$, where $B$ is the magnetic field, $\mu$ is the magnitude of the magnetic moment. In what follows we will assume that $\mu = 1$ and call $h$ the magnetic field.

The magnetization of the system $m$ will be called the average value of the magnetic moment (in units $\mu$): $$m = \frac{1}{N}\sum_i s_i = \langle s_i \rangle.$$ Summation is over all lattice nodes, $N = L^2$ is the number of nodes.

The magnetic susceptibility is defined as $$\chi = \left.\frac{\partial m}{\partial h}\right|_{h=0}.$$

The correlation function of magnetic moments is the average of the product of two magnetic moments $s_i$ and $s_j$ shifted by a vector $\vec{l}$ relative to each other (taking into account that nodes on opposite edges of the lattice are identified): $$C(\vec{l}) = \frac{1}{N}\sum_i s_i s_j.$$ This value depends on the displacement vector and on the system parameters (size, temperature, magnetic field).

\begin{figure}[H]
    \centering
    \includegraphics[width=0.72\textwidth]{PE2_EN_assets/fig1.jpg}
    \caption{Figure 1.}
\end{figure}

We will work with a lattice of size $L\times L$. The total energy of the system is $$E = -J \sum_{\langle ij \rangle} s_i s_j - h \sum_i s_i;$$ the summation is over all pairs of neighboring magnetic moments, with the assumption that magnetic moments located on opposite edges of the lattice interact with each other (periodic boundary conditions). Thanks to this, the boundary nodes are not distinguished. For example, for the lattice shown in the figure, for node 5 (located in the center) the neighbors are nodes 2, 4, 6, 8. For node 2 the neighbors are 1, 3, 5, 8. Finally, for node 3 the neighbors are 2, 6, 1, 9. The second term takes into account the external magnetic field $h$, with summation over all lattice sites.

\begin{figure}[H]
    \centering
    \includegraphics[width=0.72\textwidth]{PE2_EN_assets/fig2.jpg}
    \caption{Figure 2.}
\end{figure}

The energy of the system is minimal when all magnetic moments are oriented in the same way - either all up or all down. However, if the temperature of the system is non-zero, it can be in a state with more energy.  The probability that the system is in a certain state is proportional to $$P \sim e^{-E/k_B T}.$$ In what follows we will set $k_B = 1$. We also set $J = 1$, so the temperature becomes a dimensionless quantity. Despite the fact that the probability of being in states with higher energy is small, there are very many such states. Therefore, at sufficiently high temperatures, the system will almost certainly be in a state in which the average magnetization is zero. At a certain temperature $T_c$ a ferromagnetic-paramagnetic phase transition will occur.

Due to the huge number of states ($2^{N}$!), it is almost impossible to determine the average values using the sum over all distributions. Therefore, we will use an approximate method - the Metropolis algorithm. The essence of the algorithm is as follows. We start with a random distribution of magnetic moments. Next, we move along the lattice and for each magnetic moment we determine the change in energy $\Delta U$ if we change the direction to the opposite. If $\Delta U< 0$, we flip the magnetic moment, otherwise we flip it with probability $p = \exp(- \Delta U/T)$. It can be shown that after a large number of iterations we will obtain all configurations of magnetic moments with the correct probabilities.  The simulation results can be used to calculate various thermodynamic quantities.

\section*{Working with the program}

\noindent\textbf{exe} The executable file for Windows is published on GitHub.

The program provided to you initializes a random distribution of magnetic moments on a square lattice of a given size $L \times L$ and implements a given number of steps of the Metropolis algorithm for a given temperature and magnetic field. At each step, the algorithm tries to flip $L^2$ randomly selected magnetic moments. After finishing the simulation, the program produces the average magnetization value $m$, as well as $m^2$ and $m^4$. These averages are calculated as follows. After each step of the algorithm, the magnetization value $m_a$ is calculated ($a$ is the step number). Then the values are displayed $$ m : \frac{1}{I}\sum_{a=1}^I m_a; $$ $$ m ^2:  \frac{1}{I}\sum_{a=1}^I m_a^2; $$ $$m^4: \frac{1}{I}\sum_{a=1}^I m_a^4.$$ The program also allows you to view a graph of the magnetization $m_a$ versus the algorithm step number, and the distribution of magnetic moments at the end of the simulation.

After launch, the program will prompt you to select the required action. This can be done by entering the desired number and pressing ENTER. Number meanings:

\begin{itemize}
    \item[\textbf{1}] -- Set the lattice size $L$ and initialize all magnetic moments (each magnetic moment is set up or down with probability $1/2$). The values of temperature and magnetic field are not changed! The program will ask you to enter the lattice size -- a natural number.
    \item[\textbf{2}] -- Run the simulation. The program will ask you to enter the number of algorithm steps and execute it for the previously set parameters. After the simulation finishes, the program outputs its results. If lattice initialization (option 1) was not performed, the magnetic moments from the previous simulation are used as initial values.
    \item[\textbf{3}] -- Display the plots obtained in the last simulation. First the distribution of magnetic moments will be shown, and then (after closing that figure) the graph of magnetization versus step number.
    \item[\textbf{4}] -- Compute the correlation function $C(\vec{l})$ using the result of the last simulation (the final distribution of magnetic moments). The program will ask you, in turn, for the components $x$ and $y$ of the displacement vector $\vec{l}$ -- integers.
    \item[\textbf{5}] -- Set the magnitude of the magnetic field. The program will ask you to enter a new magnetic field value -- a real number.
    \item[\textbf{6}] -- Set the temperature. The program will ask you to enter a new temperature value -- a positive real number.
    \item[\textbf{7}] -- Show current system parameters.
\end{itemize}

In cases where the program asks you to enter a specific parameter, do this by typing the appropriate number and pressing ENTER. If the number is a fraction, separate the whole part from the fractional part using a dot (for example, 10.5). After completing the current operation, the program will prompt you to enter the number of the next action.

\section*{Part A: Observations (2 points)}

\problembox{A0}{Determine the grating size $L$, temperature $T$ and magnetic field $h$ given as default parameters in the program.}{0.20}

\problembox{A1}{How does the execution time of program $\tau$ depend on the size of the mesh $L$ and on the number of repetitions of the algorithm $I$? How many $I = 500$ steps of the algorithm were executed on a lattice of size $L = 1000$?}{0.40}

\problembox{A2}{Observe how the distribution of magnetic moments changes with temperature. Estimate the temperature $T_c$ at which the phase transition occurs.}{0.50}

\problembox{A3}{What number of steps of the algorithm $I$ is required to obtain an equilibrium distribution at $T= 1.5$, $h=0.1$ from the initial random distribution? Default grid size.}{0.20}

\problembox{A4}{In this problem you will have to work a lot with random variables. To get reasonable answers, it is important to properly estimate random errors. As an example, consider a mesh of size $L = 5$ at temperature $T = 3.66$. How does the error in determining $m^2$ depend on the number of steps of the algorithm $I$? Determine $m^2$ with an error no greater than $0.01$.}{0.70}

\section*{Part B. Measurements (8.5 points)}

Due to the small size of the system, all measured quantities fluctuate greatly, especially near the phase transition point.  In order to obtain the parameters we are interested in with reasonable accuracy, it is necessary to make a correction for the final size of the system. Do not use meshes larger than $L = 20$ in this section!  Working with them will take too much time. \textbf{For all measurements, indicate the lattice size at which they were performed!}

The Binder parameter is the quantity $$U = 1 - \frac{\langle m^4 \rangle}{3 \langle m^2 \rangle^2}.$$ The coefficients are chosen so that the parameter is zero for Gaussian fluctuations. It can be shown that if one plots the dependence of $U$ on temperature for different sizes of the system, then the curves intersect at one point exactly at the temperature of the phase transition.

\problembox{B1}{Find the value $U_0$ of the Binder parameter at $T \to 0$ and the value $U_\infty$ at $T \to \infty$.}{0.50}

\problembox{B2}{How does the mean-squared magnetization $m^2$ depend on temperature at $L = 10$? Plot a graph as a function of temperature. Select the temperature range so that the graph contains all the characteristic features of the function.}{1.00}

\problembox{B3}{Using the Binder parameter, determine as accurately as possible the temperature $T_c$ at which the phase transition occurs. Estimate the error of the found value $T_c$.}{4.00}

\problembox{B4}{Let the system be in the paramagnetic phase (the temperature is greater than the phase transition temperature $T_c$). Investigate how the magnetic susceptibility $\chi$ depends on temperature. Propose a type of dependence and determine its parameters.}{3.00}

\section*{Part C. Correlation functions (2.5 points)}

When the distance between magnetic moments is large enough, their correlation function behaves as $$C(\vec{l}) \sim \frac{1}{r^a}\, e^{-|\vec{l}|/\xi}.$$ Here $\xi$ is the correlation length. In principle it may depend on the temperature. The exponential changes much faster than the power-law factor, so to good accuracy the dependence can be considered purely exponential.

Choose the temperature as close as possible to the phase transition temperature $T = T_c$. If you were unable to determine the phase transition temperature in Part B, use the best estimate you obtained as the temperature. Explicitly state the temperature value used.

\problembox{C1}{Determine the value of the correlation length for the direction along the side of the square lattice. Over what range of distances does the exponential relationship apply?}{2.50}

\vspace{1em}
\noindent\textit{Source:} \url{https://pho.rs/p/243}

\end{document}